% ==========================================================
% titles/title-04.tex -- Title IV: Accountability Infrastructure
% ==========================================================

\ConstTitle{Accountability Infrastructure}{title:accountability}

\Article{Citizen Oversight Panels}{art:citizen-oversight-panels}

\ArtSection{Composition and Scope}{art:citizen-oversight-panels:s-composition-scope}
\gls{gls:citizen-oversight-panels}s are embedded within every governance body and \gls{gls:implementation-councils}. Each Panel consists of twelve persons selected by sortition, rotating annually, with no other governance role. Panels have full access to all internal deliberations, communications, and decisions of the body they oversee.

\ArtSection{Powers}{art:citizen-oversight-panels:s-powers}
Panels may report publicly on compliance with constitutional mandates, trigger investigation through the \gls{gls:whistleblower-infrastructure} under Article \ref{art:transparency-engine}, \ref{art:transparency-engine:s-whistleblower-infrastructure}, petition the \gls{gls:constitutional-court} for extraordinary review, and issue public findings distributed by the \gls{gls:transparency-engine} to all residents.

\ArtSection{Limits}{art:citizen-oversight-panels:s-limits}
Panels may not make binding governance decisions. Their power is communicative and investigative, not directive.

\ArtSection{Cross-Panel Coordination}{art:citizen-oversight-panels:s-cross-panel-coordination}
A \gls{gls:cross-panel-coordination-body}, composed of one rotating representative from each active \gls{gls:citizen-oversight-panels}, meets quarterly to share observed patterns across governance bodies. It has no independent authority and makes no binding decisions. All its meetings and communications are fully transparent through the \gls{gls:transparency-engine}.

\Article{The Transparency Engine}{art:transparency-engine}

\ArtSection{Definition}{art:transparency-engine:s-definition}
The \gls{gls:transparency-engine} is an open-source, publicly auditable, algorithmically automated, and cryptographically tamper-evident system through which all exercises of power at every governance level are logged, officiated, and disclosed to the public in real time. It operates without human discretion over what to disclose; disclosure is automatic and universal subject only to the narrow exceptions defined in this Article.

The act of logging itself is what makes the exercise of power legally binding.

\ArtSection{Significance Thresholds for Logging}{art:transparency-engine:s-significance-thresholds-logging}
The Constitution mandates the following logging categories that apply to all governance bodies without exception:
\begin{itemize}
  \item All governance decisions (since logging makes them official)
  \item All resource allocations
  \item All communications regarding pending or active governance decisions
  \item All deviations from pre-approved protocols
\end{itemize}

The \gls{gls:transparency-engine-oversight-body} -- a standard body of nine persons serving two-year renewable terms subject to the reconfirmation mechanism of Article \ref{art:governance-principles} -- may add logging categories but may not remove constitutionally mandated ones.

\ArtSection{Public Alert Function}{art:transparency-engine:s-public-alert-function}
The \gls{gls:transparency-engine} issues immediate public alerts to all residents whenever any purported dissolution of the constitutional order is attempted outside established processes (see Article \ref{art:antientrench}, \ref{art:antientrench:s-dissolution-prohibition}), the Anomaly Detection System identifies accumulation patterns above defined thresholds, any governance body attempts to act outside its defined authority, or emergency protocols are activated.

\ArtSection{Anomaly Detection}{art:transparency-engine:s-anomaly-detection}
The \gls{gls:transparency-engine} includes an Anomaly Detection System that monitors all resource flows for accumulation patterns inconsistent with foundational principles, all decisions for systematic benefit to identifiable groups, and all governance communications for coordination inconsistent with independent deliberation. Reports are automatic and public. The system cannot be silenced by any human actor.

\ArtSection{Whistleblower Protection}{art:transparency-engine:s-whistleblower-protection}
Any person who identifies violations of foundational principles has a protected, anonymous channel to trigger investigation through the \gls{gls:whistleblower-infrastructure} defined in \ref{art:transparency-engine:s-whistleblower-infrastructure} of this Article. Whistleblower protection is constitutionally guaranteed and cannot be overridden by any governance layer; retaliation is addressed in \ref{art:transparency-engine:s-whistleblower-infrastructure} of this Article.

\ArtSection{The Whistleblower Infrastructure}{art:transparency-engine:s-whistleblower-infrastructure}
The \gls{gls:whistleblower-infrastructure} is a constitutionally protected anonymous reporting channel embedded within the \gls{gls:transparency-engine}'s technical architecture, governed independently by a dedicated \gls{gls:whistleblower-infrastructure-panel}. Any person who identifies a violation of foundational principles may submit a report through the Infrastructure without revealing their identity to any governance body, court, or other authority.

The Infrastructure uses zero-knowledge proof of civic membership as its anonymity mechanism: a reporter may cryptographically demonstrate they are a civic member without disclosing who they are. The specific technical implementation of this mechanism is determined by the \gls{gls:whistleblower-infrastructure-panel} and published in full through the \gls{gls:transparency-engine}. The anonymity guarantee is absolute and may not be overridden by any governance body, court order, emergency declaration, or legislative majority.

The \gls{gls:whistleblower-infrastructure-panel} is a standard body of five members serving three-year terms, selected from the \gls{gls:qualification-commons} pool for security investigation competency. The Panel governs the Infrastructure, sets its operational priorities, and is responsible for its constitutional integrity. A professional technical staff implements and maintains the Infrastructure's architecture under a Staff Director selected through the \gls{gls:qualification-commons}, serving five-year terms renewable subject to the reconfirmation mechanism of Article \ref{art:governance-principles}. The Staff Director has operational independence within Panel-authorized parameters; any Panel attempt to direct specific technical implementation rather than governance priorities is automatically logged through the \gls{gls:transparency-engine} as a potential constitutional violation. The Panel coordinates with the \gls{gls:transparency-engine-oversight-body} on shared technical substrate. Neither body may access the other's exclusive data without \gls{gls:constitutional-court} authorization.

Reports submitted through the Infrastructure are routed simultaneously to the \gls{gls:citizen-oversight-panels} of the governance body identified in the report, and to the \gls{gls:security-investigation-governing-board} under Article \ref{art:internal-security}, \ref{art:internal-security:s-investigation-body} for independent assessment of whether an investigation should be opened. Neither the body identified in the report nor any person within it receives any information about the report's existence or content during any active investigation period.

The identity of a whistleblower may never be compelled, disclosed, or inferred by any governance body, court, investigation, or emergency authority at any time, for any reason. Retaliation against a whistleblower is a constitutional violation prosecuted as a matter of priority under Article \ref{art:internal-security}. Attempted identification of a whistleblower through traffic analysis, behavioral monitoring, or any other method is treated as retaliation regardless of whether identification is achieved.

\Article{The Civic Initiative Infrastructure}{art:civic-initiative-infrastructure}

\ArtSection{Definition and Purpose}{art:civic-initiative-infrastructure:s-definition-purpose}
The \gls{gls:civic-initiative-infrastructure} is an open-source, publicly auditable Commons institution that provides the civic membership with verified, manipulation-resistant channels to formally trigger defined constitutional interventions. It operates as the receive counterpart to the \gls{gls:transparency-engine}'s (Article \ref{art:transparency-engine}) broadcast function, using shared technical substrate under distinct governance. It is not a general petition or referendum platform; its scope is limited to the constitutional intervention types defined in this Article.

\ArtSection{Governance}{art:civic-initiative-infrastructure:s-governance}
The Infrastructure is governed by a nine-member \gls{gls:civic-initiative-oversight-body}, selected by sortition, serving staggered three-year terms renewable subject to the reconfirmation mechanism of Article \ref{art:governance-principles}. The Oversight Body maintains the Infrastructure's technical integrity, publishes full documentation of all processes, and may not alter intervention thresholds, cooling periods, or triggered consequences, which are constitutionally fixed. All deliberations and decisions of the Oversight Body are logged in real time on the \gls{gls:transparency-engine}.

\ArtSection{Standard Mechanics}{art:civic-initiative-infrastructure:s-standard-mechanics}
Any civic member may initiate a formal intervention petition through the Infrastructure. Upon initiation, the petition is immediately logged publicly through the \gls{gls:transparency-engine}, displaying its existence, target, stated grounds, and running verified signature count in real time. Signatures are verified against the civic membership roll by the \gls{gls:electoral-commons} under Article \ref{art:sortition-design}, \ref{art:sortition-design:s-electoral-commons} as submitted. Individual signatories are disclosed only to the \gls{gls:electoral-commons} during the active petition period; full signatory lists become public through the \gls{gls:transparency-engine} upon expiration of the relevant cooling period. Anonymous signatures are not permitted. The Infrastructure automatically enforces all cooling periods; no governance body or petitioner group may waive them. Upon threshold confirmation, the Infrastructure automatically notifies the relevant governance body, the relevant \gls{gls:citizen-oversight-panels}, and the \gls{gls:constitutional-court}. The triggered process commences within thirty days of threshold confirmation.

\ArtSection{Constitutional Intervention Types}{art:civic-initiative-infrastructure:s-constitutional-intervention-types}
The following intervention types are constitutionally established. The \gls{gls:sortition-legislature} may add additional intervention types, but may not remove the ones listed below. Individual articles specify any domain-specific constraints.

\textbf{\gls{gls:constitutional-court} Decision Review}
\begin{itemize}
  \item Requires five percent of civic membership
  \item Applicable once per Court decision
  \item Triggers constitution of a new sortition Court with no membership overlap from the original body, which will examine the previous Court's decision
\end{itemize}

\textbf{Technical Position Review}
\begin{itemize}
  \item Requires three percent of civic membership
  \item Subject to an eighteen-month cooling period per position
  \item Triggers full reconfirmation process under Article \ref{art:governance-principles}
\end{itemize}
This intervention type applies to positions designated as renewable under Article \ref{art:governance-principles}. Non-renewable position holders are subject to removal only through the processes established in Article \ref{art:sortition-legislature}, \ref{art:sortition-legislature:s-bad-faith-removal} and Article \ref{art:const-court} as applicable.

\textbf{Legislative Decision Review}
\begin{itemize}
  \item Requires five percent of civic membership
  \item Subject to a twenty-four-month cooling period per decision
  \item Triggers a mandatory legislative reconsideration session within sixty days
  \item Reconsideration does not suspend the original decision and does not bind the Legislature to reversal
\end{itemize}

\textbf{Regulatory Body Mandate Review}
\begin{itemize}
  \item Requires four percent of civic membership
  \item Subject to a twenty-four-month cooling period per Body
  \item Triggers \gls{gls:constitutional-court} review of whether the Body has acted within its legislatively defined mandate
\end{itemize}

\textbf{Emergency Protocol Post-hoc Review}
\begin{itemize}
  \item Requires three percent of civic membership
  \item Applicable once per emergency activation
  \item Triggers constitution of a one-time \gls{gls:emergency-review-panel}
\end{itemize}

The \gls{gls:emergency-review-panel} consists of one hundred fifty persons selected by sortition from the civic membership, with no membership overlap with the Legislature that authorized the emergency actions under review. The Panel convenes within thirty days of threshold confirmation, reviews all actions taken under the relevant emergency authorization for compliance with constitutional mandates, and submits findings to the \gls{gls:constitutional-court} (Article \ref{art:const-court}) within ninety days. The \gls{gls:constitutional-court} issues binding determinations on the constitutionality of reviewed actions within sixty days of receiving findings. The Panel dissolves upon submission of findings. The existing Legislature continues its normal operations throughout this process.

\ArtSection{Protections}{art:civic-initiative-infrastructure:s-protections}
Retaliation against any person for initiating, signing, or publicly supporting a civic intervention petition is a constitutional violation prosecuted as a matter of priority under Article \ref{art:internal-security}. The Infrastructure's technical architecture is designed to prevent identification of signatories by any governance body during the active petition period. Any attempt by any governance body to access signatory data outside the \gls{gls:electoral-commons} verification process is automatically logged by the \gls{gls:transparency-engine} and referred to the \gls{gls:security-enforcement-review-panel}.

\ArtSection{Relationship to the Transparency Engine}{art:civic-initiative-infrastructure:s-relationship-transparency-engine}
The \gls{gls:civic-initiative-infrastructure} and the \gls{gls:transparency-engine} share technical infrastructure and are jointly maintained by their respective Oversight Bodies under a published cooperation agreement. Neither Body may access the other's exclusive data without authorization from the \gls{gls:constitutional-court}. The cooperation agreement is fully public and may not be amended without supermajority approval from both Oversight Bodies and \gls{gls:constitutional-court} review.